\makeatletter
\def\input@path{{./}{../}{../../}{../../../}{../../../../}{preamble/}{../preamble/}{../../preamble/}{../../../preamble/}}
\makeatother

\input{preamble_exams.tex}
\input{graphs/highlights.tex}
% \GTenHighlightsfalse % Uncomment to disable instructional highlights.
\input{graphs/specific_highlights.tex}
\CoordinatedTermTwoPalette
\usepackage{needspace}
\usepackage{titling}
\setlength{\droptitle}{-2em}
\pretitle{\par\begingroup\centering\LARGE}
\posttitle{\par\endgroup}
\preauthor{}\postauthor{}
\predate{}\postdate{}
\maketitlehookd{\vspace{-1em}}
\linespread{1.06}
\setlength{\parskip}{4pt}
\newgeometry{margin=0.65in,top=0.55in,bottom=0.65in}

\title{Grade 11 Term 2 Home Practice -- C5\\\large A population proportion\\[0.25em]
	\normalsize Global Economics -- 2026--2027}
\author{}
\date{}


% Local refinement of the existing registry entry: essential interval roles only.
% Shared registry and coordinated palette remain unchanged for other activities.
\expandafter\renewcommand\csname practicehighlightlegend@G11-T2-C5-coordinated\endcsname{%
  \begingroup
  \expandafter\renewcommand\csname coordinatednotation@11@se\endcsname{$\HLStandardError{SE}$}%
  \expandafter\renewcommand\csname coordinatednotation@11@critical\endcsname{$\HLCriticalValue{z^*}$}%
  \expandafter\renewcommand\csname coordinatednotation@11@margin\endcsname{$\HLMarginOfError{E}$}%
  \expandafter\renewcommand\csname coordinatednotation@11@propsample\endcsname{$\HLSampleProportion{\hat p}$}%
  \CoordinatedLegendTable{%
    $\HLPopulationProportion{p}$ & \HLPopulationProportion{Population} \HLPopulationProportion{proportion} & Unknown population fraction.\\
    $\HLSampleProportion{\hat p}$ & \HLSampleProportion{Sample} \HLSampleProportion{proportion} & Observed sample fraction.\\
    $\HLObservationCount{x},\ \CtwoSize{n},\ \HLPopulationGroup{N}$ & \HLObservationCount{Sample counts}; \HLPopulationGroup{population size} & Successes; sample size; population size.\\
    $\HLStandardError{SE}$ & \HLStandardError{Standard error} & Estimated sampling spread.\\
    $\HLConfidence{C}$ & \HLConfidence{Confidence level} & Approximate long-run coverage.\\
    $\HLCriticalValue{z^*}$ & \HLCriticalValue{Critical value} & Standard normal multiplier.\\
    $\HLMarginOfError{E}$ & \HLMarginOfError{Margin of error} & Distance from centre to endpoint.\\
    $\HLIntervalBounds{L},\ \HLIntervalBounds{U}$ & \HLIntervalBounds{Endpoints} & Lower and upper bounds for $\HLPopulationProportion{p}$.\\
    $\HLCumulativeProbability{\Phi(z)}$ & \HLCumulativeProbability{Cumulative area} & Standard normal area to the left of $z$.\\
    $\HLIntervalWidth{W}$ & \HLIntervalWidth{Full width} & Distance between the endpoints.\\
  }\endgroup}
\newcommand{\IntroOne}{\hyperref[ex:intro-one]{Example 1}}
\newcommand{\IntroTwo}{\hyperref[ex:intro-two]{Example 2}}

\begin{document}
\maketitle
\setcounter{tocdepth}{1}
\setcounter{secnumdepth}{2}
\phantomsection\label{toc}
\tableofcontents

\section*{Evaluated criteria}
\addcontentsline{toc}{section}{Evaluated criteria}
\begin{enumerate}[label=\textbf{C\arabic*},start=5]
\item Constructs confidence intervals for a population proportion using sample proportions, applying them to contexts such as poverty rates, defaulting borrowers, or successful investments.
\end{enumerate}
\textbf{How to practise.} Try each task before opening its blue solution link. Show the model, conditions, formula, algebra, intermediate values and a contextual conclusion. Use decimals in formulas, retain full precision until the final display, and report confidence as a percentage. The arrows return to the contents, problem list or current problem.

This activity uses the symmetric, unadjusted one-proportion normal ($z$) interval. The inverse problems recover the inputs to this same procedure. All problems address C5.
\clearpage

\section{Introduction}\label{sec:introduction}
\subsection*{Highlight legend}
\phantomsection\label{tab:highlight-legend}
\PracticeHighlightLegend{G11-T2-C5-coordinated}
Names and notation retain meaning without color. Related shades connect related quantities: confidence and a critical value are connected but have different roles.

\subsection{Introductory example 1: borrowers and defaults}\label{ex:intro-one}
A credit cooperative randomly selects \CtwoSize{200} loans without replacement from \HLPopulationGroup{4,000} loans with complete annual records. It observes \HLObservationCount{30} first-year defaults. Here a ``success'' means a default, even though that is economically undesirable. The \HLSampleProportion{sample proportion} is
\[
\HLSampleProportion{\hat p}=\frac{\HLObservationCount{x}}{\CtwoSize{n}}
=\frac{\HLObservationCount{30}}{\CtwoSize{200}}=\HLSampleProportion{0.15}.
\]
The unknown \HLPopulationProportion{population proportion} $\HLPopulationProportion{p}$ describes all 4,000 loans. A count of 30 loans is different from a proportion of 0.15, or 15\%. Sampling 5\% of the population supports approximate independence; the observed counts are 30 and 170, both at least 10.

\subsection{Introductory example 2: interpreting shipment uncertainty}\label{ex:intro-two}
An export council randomly selects \CtwoSize{400} shipments without replacement from \HLPopulationGroup{10,000}; \HLObservationCount{80} arrive late. Thus $\HLSampleProportion{\hat p}=\HLSampleProportion{0.20}$. The 4\% sampling fraction and counts 80 and 320 support the approximation. A \HLConfidence{95\%} interval using $\HLCriticalValue{z^*}=\HLCriticalValue{1.96}$ is $[0.1608,0.2392]$, rounded. Its \HLMarginOfError{margin} is about \HLMarginOfError{3.92 percentage points}, not 3.92\% of the estimate.

\textbf{What confidence means.} If we repeatedly sampled in the same way and constructed intervals by this procedure, approximately \HLConfidence{95\%} would contain the fixed population delay proportion. An individual computed interval either contains $p$ or misses it; 95\% is the procedure's approximate long-run coverage, not a probability assigned to the fixed $p$ after observing the data.
\clearpage

\section{Notation Reference and Practice Procedure}\label{sec:notation-and-procedure}
\subsection{Notation reference}\label{tab:notation-reference}
\hyperref[tab:highlight-legend]{Return to the concise legend.} Here $SE$ always denotes an \emph{estimated} standard error, also written $\widehat{SE}$.
\begingroup
\setlength{\tabcolsep}{4pt}
\tikzset{gten marker/.append style={inner xsep=.6mm,inner ysep=.3mm}}
\renewcommand{\arraystretch}{1.25}
\begin{longtable}{@{}>{\raggedright\arraybackslash}p{.16\linewidth}>{\raggedright\arraybackslash}p{.19\linewidth}>{\raggedright\arraybackslash}p{.29\linewidth}>{\raggedright\arraybackslash}p{\dimexpr.36\linewidth-6\tabcolsep\relax}@{}}
\toprule
\textbf{Symbol or notation}&\textbf{Concept}&\textbf{Definition}&\textbf{Meaning in the introductory examples}\\\midrule
\endfirsthead
\toprule
\textbf{Symbol or notation}&\textbf{Concept}&\textbf{Definition}&\textbf{Meaning in the introductory examples}\\\midrule
\endhead
\bottomrule\endlastfoot
$\HLPopulationProportion{p}$ & \HLPopulationProportion{Population} \HLPopulationProportion{proportion} & Fixed unknown fraction of all population units with the defined outcome; $0\le p\le1$. & \textbullet\ \IntroOne: fraction of all loans defaulting.\par\textbullet\ \IntroTwo: fraction of all shipments late.\\
$\HLSampleProportion{\hat p}$ & \HLSampleProportion{Sample proportion} & Observed fraction $x/n$; the interval centre. & \textbullet\ \IntroOne: \HLSampleProportion{0.15}.\par\textbullet\ \IntroTwo: \HLSampleProportion{0.20}.\\
$\HLObservationCount{x},\ \CtwoSize{n}$ & \HLObservationCount{Success count}; \CtwoSize{sample size} & Integer numbers of successes and sampled units; failures $n-x$. & \textbullet\ \IntroOne: \HLObservationCount{30} defaults in \CtwoSize{200} loans.\par\textbullet\ \IntroTwo: \HLObservationCount{80} late in \CtwoSize{400}.\\
$\HLPopulationGroup{N}$ & \HLPopulationGroup{Population size} & Number of units available for sampling; check $n\le0.10N$ without replacement. & \textbullet\ \IntroOne: \HLPopulationGroup{4,000} loans.\par\textbullet\ \IntroTwo: \HLPopulationGroup{10,000} shipments.\\
$\HLStandardError{SE}$ & \HLStandardError{Standard error} & Estimated sampling spread of $\hat p$: $\sqrt{\hat p(1-\hat p)/n}$. & \textbullet\ \IntroTwo: about \HLStandardError{0.0200} as a proportion.\\
$\HLConfidence{C}$ & \HLConfidence{Confidence level} & Nominal central normal area and approximate long-run coverage; decimal in formulas. & \textbullet\ \IntroTwo: \HLConfidence{0.95}, reported as \HLConfidence{95\%}.\\
$\HLCriticalValue{z^*}$ & \HLCriticalValue{Critical value} & Positive normal multiplier; $\Phi^{-1}((1+C)/2)$. & \textbullet\ \IntroTwo: \HLCriticalValue{1.96} for \HLConfidence{95\%}.\\
$\HLCumulativeProbability{\Phi(z)}$ & \HLCumulativeProbability{Cumulative area} & $P(Z\le z)$ for a standard normal variable $Z$. & \textbullet\ \IntroTwo: $\HLCumulativeProbability{\Phi(1.96)}\approx\HLCumulativeProbability{0.9750}$.\\
$\HLMarginOfError{E}$ & \HLMarginOfError{Margin of error} & Half-width $z^*SE=(U-L)/2$; measured in proportion units. & \textbullet\ \IntroTwo: \HLMarginOfError{0.0392}, or \HLMarginOfError{3.92} percentage points.\\
$\HLIntervalBounds{L},\ \HLIntervalBounds{U}$ & \HLIntervalBounds{Endpoints} & Lower and upper bounds $\hat p-E$ and $\hat p+E$. & \textbullet\ \IntroTwo: about \HLIntervalBounds{0.1608} and \HLIntervalBounds{0.2392}.\\
$\HLIntervalWidth{W}$ & \HLIntervalWidth{Full width} & $U-L=2E$; twice the margin. & \textbullet\ \IntroTwo: \HLIntervalWidth{0.0784}, or \HLIntervalWidth{7.84} percentage points.\\
\end{longtable}
\endgroup
\clearpage

\subsection{Formula map: choose the missing quantities}
All three types use the same symmetric, unadjusted one-proportion $z$ procedure, with no continuity correction or finite-population correction. Its estimated sampling spread comes from the observed proportion, not from the unknown $p$.
\begin{align*}
\HLSampleProportion{\hat p}&=\frac{\HLObservationCount{x}}{\CtwoSize{n}},&
\HLStandardError{SE}&=\sqrt{\frac{\HLSampleProportion{\hat p}(1-\HLSampleProportion{\hat p})}{\CtwoSize{n}}},\\
\HLCriticalValue{z^*}&=\Phi^{-1}\!\left(\frac{1+\HLConfidence{C}}{2}\right),&
\HLMarginOfError{E}&=\HLCriticalValue{z^*}\,\HLStandardError{SE},\\
[\HLIntervalBounds{L},\HLIntervalBounds{U}]&=[\HLSampleProportion{\hat p}-\HLMarginOfError{E},\HLSampleProportion{\hat p}+\HLMarginOfError{E}],&
\HLIntervalWidth{W}&=2\HLMarginOfError{E}.
\end{align*}
For either inverse type, adding and subtracting $L=\hat p-E$ and $U=\hat p+E$ gives
\[
\HLSampleProportion{\hat p}=\frac{\HLIntervalBounds{L}+\HLIntervalBounds{U}}{2},\qquad
\HLMarginOfError{E}=\frac{\HLIntervalBounds{U}-\HLIntervalBounds{L}}{2}.
\]
Squaring $E=z^*\sqrt{\hat p(1-\hat p)/n}$ and rearranging gives
\[
nE^2=(z^*)^2\hat p(1-\hat p)
\quad\Longrightarrow\quad
\CtwoSize{n}=\left(\frac{\HLCriticalValue{z^*}}{\HLMarginOfError{E}}\right)^2\hat p(1-\hat p).
\]
Alternatively, divide by the estimated standard error and recover confidence:
\[
\HLCriticalValue{z^*}=\frac{\HLMarginOfError{E}}{\sqrt{\hat p(1-\hat p)/n}},\qquad
\HLConfidence{C}=2\HLCumulativeProbability{\Phi(z^*)}-1.
\]
\textbf{Calculator instructions.} To find $\Phi(z^*)$, use a normal cumulative-distribution command with lower bound $-\infty$ (or $-10$), upper bound $z^*$, mean $0$, and standard deviation $1$. Use $\operatorname{invNorm}((1+C)/2,0,1)$ for the critical value. Do not use the normal density command. For checks, $\Phi(1.645)\approx0.950015$, $\Phi(1.96)\approx0.975002$ and $\Phi(2.326)\approx0.989991$. Supplied multipliers are rounded normal quantiles.

\subsection{Reusable practice procedure}
\begin{enumerate}
\item Define a success, the population and its fixed proportion $p$. Identify givens and unknowns: construct $[L,U]$; recover $\hat p,n$; or recover $\hat p,C$.
\item Check a random sample, approximate independence, and observed counts $x\ge10$ and $n-x\ge10$. Without replacement, check $n/N\le0.10$. For inverse tasks, complete these checks after recovering the counts.
\item Use the appropriate formula map. Show rearrangements, substitute decimals, and retain exact expressions or unrounded intermediate values. Confirm $0\le L<U\le1$ for these tasks.
\item Check recovered $n$ and $x=n\hat p$ are valid integers. Rounded endpoints generally allow only approximate recovery, with no unique exact sample size or count guaranteed. Do not automatically round $n$ upward: reconstruction is not minimum-sample-size planning.
\item Report proportions as percentages and margins or widths in percentage points. Interpret approximate long-run coverage and compare widths only after stating which inputs are held fixed.
\end{enumerate}
\clearpage

\section{General Worked Problems}\label{sec:worked-solution-examples}
\SectionProblemsTableOfContents{
\SectionProblemLink{prob:c5-1}{Type 1: construct a confidence interval}\\
\SectionProblemLink{prob:c5-2}{Type 2: recover the sample proportion and sample size}\\
\SectionProblemLink{prob:c5-3}{Type 3: recover the sample proportion and confidence level}
}

\subsection{Type 1: construct a confidence interval}\label{prob:c5-1}
A lender takes a simple random sample without replacement of 400 loans from 8,000 small-business loans issued in the same year. All sampled loans have complete first-year records; 50 defaulted. Use $C=0.98$ and the supplied $z^*=2.326$.
\QuestionsTasks[sec:worked-solution-examples]{
{5}/{sol:c5-1-a}/{Define $p$, calculate $\hat p$, and check the conditions.},
{5}/{sol:c5-1-b}/{Construct and interpret the 98\% interval. Compare it with a 20\% benchmark. With $\hat p$ fixed, explain the effects of quadrupling $n$ at the same confidence and of lowering confidence to 95\% at the original $n$ (use $z^*=1.96$).}
}
\subsubsection*{C5 -- Identify the parameter and check the procedure}\label{sol:c5-1-a}
The \HLPopulationProportion{parameter} $\HLPopulationProportion{p}$ is the fraction of all 8,000 loans defaulting within one year; the sample provides an estimate:
\[
\HLSampleProportion{\hat p}=\frac{\HLObservationCount{x}}{\CtwoSize{n}}=
\frac{\HLObservationCount{50}}{\CtwoSize{400}}=\HLSampleProportion{0.125}.
\]
Random selection supports inference to this cohort. The sampling fraction is $400/8{,}000=5\%\le10\%$, supporting approximate independence. Observed defaults $x=50$ and non-defaults $n-x=350$ both exceed 10.
\secInicio[prob:c5-1][sec:worked-solution-examples]

\subsubsection*{C5 -- Build, interpret and compare}\label{sol:c5-1-b}
The binary outcome and the checked conditions support the symmetric normal approximation. Using $z^*=\Phi^{-1}((1+C)/2)\approx2.326$,
\begin{align*}
\HLStandardError{SE}&=\sqrt{\frac{\hat p(1-\hat p)}{n}}
=\sqrt{\frac{\HLSampleProportion{0.125}(1-\HLSampleProportion{0.125})}{\CtwoSize{400}}}
\approx\HLStandardError{0.01653595},\\
\HLMarginOfError{E}&=\HLCriticalValue{z^*}\,\HLStandardError{SE}
=\HLCriticalValue{2.326}(0.01653594569\ldots)\approx\HLMarginOfError{0.03846261}.
\end{align*}
Then $[L,U]=[\hat p-E,\hat p+E]\approx[\HLIntervalBounds{0.08654},\HLIntervalBounds{0.16346}]$, within $[0,1]$. We are 98\% confident that about 8.65\% to 16.35\% of this cohort defaulted. Across repeated samples, this procedure would cover the fixed cohort proportion approximately 98\% of the time. The margin is about 3.85 percentage points. The entire interval is below 20\%, providing evidence that $p<0.20$.

\textbf{Follow-ups.} Since $W=2z^*\sqrt{\hat p(1-\hat p)/n}$, holding $\hat p$ and $C$ fixed and changing $n$ from 400 to 1,600 multiplies the width by $\sqrt{400/1600}=1/2$; $E\approx0.01923130$. This comparison needs a population of at least 16,000 if sampling without replacement; for example, $N=32,000$ gives 5\%. Holding $\hat p$ and $n=400$ fixed, reducing $C$ to 95\% reduces $z^*$ to 1.96 and $E$ to $1.96(0.01653594569\ldots)\approx0.03241045$. The narrower interval has lower long-run coverage.
\secInicio[prob:c5-1][sec:worked-solution-examples]
\clearpage

\subsection{Type 2: recover the sample proportion and sample size}\label{prob:c5-2}
An export council used a simple random sample without replacement from 12,000 shipments. ``Success'' means arrival at least one day late. Its interval was constructed using the \textbf{symmetric, unadjusted one-proportion $z$ procedure}, with $C=0.95$ and the supplied $z^*=1.96$. The exact endpoints are
\[
L=\frac{3}{10}-\frac{49\sqrt{14}}{5000},\qquad
U=\frac{3}{10}+\frac{49\sqrt{14}}{5000},
\]
with rounded display $[0.26333,0.33667]$.
\QuestionsTasks[sec:worked-solution-examples]{
{5}/{sol:c5-2-a}/{Recover $\hat p$, $E$ and $n$ using the exact endpoints. Recover $x=n\hat p$, check integrality and the sampling conditions, and reconstruct the interval.},
{5}/{sol:c5-2-b}/{Explain what changes if only the rounded endpoints are available. Should a recovered sample size always be rounded upward?}
}
\subsubsection*{C5 -- Recover, verify and reconstruct}\label{sol:c5-2-a}
The stated procedure makes the centre equal to the sample delay proportion and the half-width equal to its margin. Adding and subtracting the endpoints gives
\[
\HLSampleProportion{\hat p}=\frac{L+U}{2}=\frac{0.6}{2}=\HLSampleProportion{0.30},\qquad
\HLMarginOfError{E}=\frac{U-L}{2}=\frac{49\sqrt{14}}{5000}\approx\HLMarginOfError{0.03666824}.
\]
From $E=z^*\sqrt{\hat p(1-\hat p)/n}$, square both sides and rearrange:
\begin{align*}
E^2&=\frac{(z^*)^2\hat p(1-\hat p)}{n}
\quad\Longrightarrow\quad n=\left(\frac{z^*}{E}\right)^2\hat p(1-\hat p),\\
\CtwoSize{n}&=\left(\frac{\HLCriticalValue{1.96}}{49\sqrt{14}/5000}\right)^2
(\HLSampleProportion{0.30})(1-\HLSampleProportion{0.30})
=\frac{10^8}{14}\cdot\frac{21}{100}=\CtwoSize{600},\\
\HLObservationCount{x}&=n\hat p=600(0.30)=\HLObservationCount{180}.
\end{align*}
Both $n=600$ and $x=180$ are integers. The observed failure count is $600-180=420$, so both counts exceed 10. Random sampling is given; $600/12{,}000=5\%$ supports approximate independence. Reconstructing,
\[
\HLStandardError{SE}=\sqrt{\frac{0.30(0.70)}{600}}=\frac{\sqrt{14}}{200}
\approx\HLStandardError{0.01870829},\qquad
E=1.96\frac{\sqrt{14}}{200}=\frac{49\sqrt{14}}{5000}.
\]
Thus $\hat p\pm E$ reproduces the exact endpoints, both within $[0,1]$. The observed delay rate is 30\%, with a margin of about 3.67 percentage points. We are 95\% confident that about 26.33\% to 33.67\% of all seasonal shipments were late, in the approximate long-run coverage sense.
\secInicio[prob:c5-2][sec:worked-solution-examples]

\subsubsection*{C5 -- Rounded reports and honest recovery}\label{sol:c5-2-b}
Using only $[0.26333,0.33667]$ gives $\hat p\approx0.30$, $E\approx0.03667$ and
\[
n\approx(1.96/0.03667)^2(0.30)(0.70)\approx599.9425,
\qquad x\approx179.9827.
\]
These are approximate reconstructions; rounded endpoints do not generally determine unique integers. Here the exact endpoints establish 600 and 180. Do not automatically round $n$ upward as in minimum-size planning. Check candidate integers against the stated method and endpoint rounding.
\secInicio[prob:c5-2][sec:worked-solution-examples]
\clearpage

\subsection{Type 3: recover the sample proportion and confidence level}\label{prob:c5-3}
A trade association draws a simple random sample without replacement of $n=500$ retailers from 10,000 registered retailers. A success is a retailer accepting a new digital payment service. Its interval was constructed using the \textbf{symmetric, unadjusted one-proportion $z$ procedure}, with exact endpoints
\[
L=\frac15-\frac{329\sqrt5}{25000},\qquad
U=\frac15+\frac{329\sqrt5}{25000}.
\]
The rounded display is $[0.17057,0.22943]$; the confidence level is not supplied.
\QuestionsTasks[sec:worked-solution-examples]{
{5}/{sol:c5-3-a}/{Recover $\hat p$ and $E$, find the estimated standard error and $z^*$, then use $\Phi(z^*)$ to calculate $C$. Verify the counts, sampling conditions and endpoints; interpret the interval.},
{5}/{sol:c5-3-b}/{Holding $n$ and $C$ fixed, compare the full width at $\hat p=0.20$ with the width at $\hat p=0.50$. Explain why the width changes.}
}
\subsubsection*{C5 -- Recover the confidence level}\label{sol:c5-3-a}
Here $p$ is the fraction of all registered retailers accepting the service. Symmetry identifies the sample fraction and margin:
\[
\HLSampleProportion{\hat p}=\frac{L+U}{2}=\frac{2/5}{2}=\HLSampleProportion{0.20},\qquad
\HLMarginOfError{E}=\frac{U-L}{2}=\frac{329\sqrt5}{25000}\approx\HLMarginOfError{0.02942665}.
\]
The observed success count is $x=n\hat p=500(0.20)=100$, an integer, with 400 failures. Both exceed 10. Sampling is random and $500/10{,}000=5\%$, so approximate independence is reasonable. Use the estimated standard error for this binary outcome:
\begin{align*}
\HLStandardError{SE}&=\sqrt{\frac{\hat p(1-\hat p)}{n}}
=\sqrt{\frac{\HLSampleProportion{0.20}(1-\HLSampleProportion{0.20})}{\CtwoSize{500}}}
=\frac{\sqrt5}{125}\approx\HLStandardError{0.01788854},\\
E&=z^*SE\quad\Longrightarrow\quad
\HLCriticalValue{z^*}=\frac{E}{SE}
=\frac{329\sqrt5/25000}{\sqrt5/125}=\HLCriticalValue{1.645}.
\end{align*}
The normal cumulative calculator gives $\HLCumulativeProbability{\Phi(1.645)}\approx\HLCumulativeProbability{0.95001509}$. Because the central area excludes two equal tails,
\[
\HLConfidence{C}=\Phi(z^*)-\Phi(-z^*)=2\Phi(z^*)-1
\approx2(0.95001509)-1=\HLConfidence{0.90003019}.
\]
Report the confidence level as approximately \HLConfidence{90.003\%}, conventionally \HLConfidence{90\%}. The small difference comes from the rounded multiplier 1.645. Multiplying $1.645(\sqrt5/125)$ recovers the exact margin and endpoints, within $[0,1]$. We are about 90\% confident that 17.06\% to 22.94\% of registered retailers accept the service. Repeated use of this sampling-and-interval procedure would cover the fixed $p$ about 90\% of the time.
\secInicio[prob:c5-3][sec:worked-solution-examples]

\subsubsection*{C5 -- How the sample proportion affects width}\label{sol:c5-3-b}
At fixed $n=500$ and $z^*=1.645$, $W=2z^*\sqrt{\hat p(1-\hat p)/n}$. Changing $\hat p$ from 0.20 to 0.50 multiplies $W$ by $\sqrt{0.25/0.16}=1.25$: from about 0.05885331 to 0.07356664, or from 5.89 to 7.36 percentage points. Since $\hat p(1-\hat p)=1/4-(\hat p-1/2)^2$, width is largest at $\hat p=0.50$ and equal at complementary proportions such as 0.20 and 0.80. At 0.50 the counts would be 250 and 250, still valid.
\secInicio[prob:c5-3][sec:worked-solution-examples]
\clearpage

\section{Advanced Problems}\label{sec:advanced}
\SectionProblemsTableOfContents{
\SectionProblemLink{prob:c5-advanced}{Compare two seasonal default-rate reports}
}
\subsection{Compare two seasonal default-rate reports}\label{prob:c5-advanced}
Two lenders each use a simple random sample without replacement from a separate cohort of 30,000 loans with complete first-year records. A success is a first-year default. Both reports use the \textbf{symmetric, unadjusted one-proportion $z$ procedure}. All endpoints below are exact for the stated multipliers.
\begin{center}
\begin{tabular}{@{}lll@{}}\toprule
Report & Interval & Additional information\\\midrule
A & $[0.201,0.299]$ & $C=0.95$, use $z^*=1.96$; $n$ unknown\\
B & $[0.220925,0.279075]$ & $n=1,200$; $C$ unknown\\\bottomrule
\end{tabular}
\end{center}
\QuestionsTasks[sec:advanced]{
{5}/{sol:c5-advanced-a}/{Recover $\hat p$, $n$ and $x$ for A. Recover $\hat p$, $x$, $SE$, $z^*$ and $C$ for B. Verify the conditions and reconstruct both intervals.},
{5}/{sol:c5-advanced-b}/{Compare the widths. Explain how B can have both higher confidence and a narrower interval. Is either report evidence that its cohort default rate is below 20\%? Interpret the confidence and compare the estimates without claiming a tested difference between populations.}
}
\subsubsection*{C5 -- Coordinate the two inverse methods}\label{sol:c5-advanced-a}
Each unknown $p$ is its lender's population default fraction. The stated model allows midpoint and half-width recovery. For A,
\begin{align*}
\HLSampleProportion{\hat p_A}&=(0.201+0.299)/2=\HLSampleProportion{0.25},&
\HLMarginOfError{E_A}&=(0.299-0.201)/2=\HLMarginOfError{0.049},\\
E_A^2&=(z_A^*)^2\hat p_A(1-\hat p_A)/n_A,&&\\
\CtwoSize{n_A}&=(z_A^*/E_A)^2\hat p_A(1-\hat p_A)\\
&=(1.96/0.049)^2(0.25)(0.75)=\CtwoSize{300},\\
\HLObservationCount{x_A}&=n_A\hat p_A=300(0.25)=\HLObservationCount{75}.&&
\end{align*}
Thus $SE_A=\sqrt{0.25(0.75)/300}=0.025$ and $z_A^*SE_A=1.96(0.025)=0.049$. The resulting $0.25\pm0.049$ reproduces A.

For B, the same midpoint and half-width formulas give
\begin{align*}
\HLSampleProportion{\hat p_B}&=(0.220925+0.279075)/2=\HLSampleProportion{0.25},\\
\HLMarginOfError{E_B}&=(0.279075-0.220925)/2=\HLMarginOfError{0.029075},\\
\HLObservationCount{x_B}&=1200(0.25)=\HLObservationCount{300},\qquad
\HLStandardError{SE_B}=\sqrt{0.25(0.75)/1200}=\HLStandardError{0.0125},\\
E_B&=z_B^*SE_B\quad\Longrightarrow\quad
\HLCriticalValue{z_B^*}=E_B/SE_B=0.029075/0.0125=\HLCriticalValue{2.326},\\
\HLConfidence{C_B}&=2\Phi(2.326)-1\approx2(0.98999072)-1
=\HLConfidence{0.97998145}\approx\HLConfidence{98\%}.
\end{align*}
Reconstruction gives $0.25\pm2.326(0.0125)=[0.220925,0.279075]$. Both intervals are within $[0,1]$. Recovered counts are integers: A has 75 defaults and 225 non-defaults; B has 300 and 900, all at least 10. Random sampling is given. Fractions $300/30{,}000=1\%$ and $1200/30{,}000=4\%$ support approximate independence.
\secInicio[prob:c5-advanced][sec:advanced]

\Needspace{18\baselineskip}
\subsubsection*{C5 -- Compare precision and conclusions}\label{sol:c5-advanced-b}
Full widths are $\HLIntervalWidth{W_A}=2(0.049)=\HLIntervalWidth{0.098}$ and $\HLIntervalWidth{W_B}=2(0.029075)=\HLIntervalWidth{0.05815}$: 9.8 and 5.815 percentage points. With $\hat p=0.25$ fixed,
\[
\frac{W_B}{W_A}=\frac{z_B^*}{z_A^*}\sqrt{\frac{n_A}{n_B}}
=\frac{2.326}{1.96}\sqrt{\frac{300}{1200}}
\approx0.59337.
\]
Holding $\hat p$ and confidence fixed, quadrupling $n$ halves the width. Holding $\hat p$ and $n$ fixed, raising confidence increases the critical value and the width by a factor $2.326/1.96\approx1.18673$. Here the sample-size effect outweighs the confidence effect, so B is about 40.66\% narrower despite its higher confidence.

Both sample estimates are 25\%; B is more precise by this width measure. We are 95\% confident that A's cohort default rate is 20.1\%--29.9\%, and about 98\% confident that B's is 22.0925\%--27.9075\%. These levels describe approximate coverage over repeated use of each procedure. Neither interval supports a default rate below 20\%: both are entirely above it. Equal sample estimates do not prove equal population rates, and these separate intervals do not constitute a formal test of the difference between the lenders.
\secInicio[prob:c5-advanced][sec:advanced]
\end{document}
